% =====================================================================
% 10. Testing, Evaluation and Preliminary Results
% =====================================================================
\section{Testing, Evaluation and Preliminary Results}
\label{sec:evaluation}

\subsection{Testing strategy}\label{subsec:testing-strategy}
Software testing and research evaluation are kept separate.
Testing checks that the artefact is correct; evaluation answers the research questions.

\begin{itemize}
    \item \textbf{Unit tests} for adapters, calibration wrappers, filters, KPI calculations and alert rules, with property-based tests for filters (for example, a constant signal must pass through unchanged).
    Target: 80\% line coverage on the processing layer.
    \item \textbf{Integration tests} run the full pipeline on a short replayed trace and compare the KPIs and alerts produced with expected values.
    \item \textbf{System tests} run long replays (several hours) to check memory stability and reconnection (NFR8).
    \item \textbf{Acceptance:} each requirement's acceptance criteria are checked at the end of the iteration that delivers it.
    Acceptance by the STT is a company activity, not research data.
\end{itemize}

\subsection{Evaluation experiments}\label{subsec:evaluation-experiments}
All experiments run on replayed or synthetic traces with seeded faults, so every result has ground truth.

\begin{description}[leftmargin=1.2em, labelindent=0pt, font=\normalfont\bfseries]
  \item[Experiment 1: Performance (RQ1).] A load generator replays configurations from a single pack at nominal rate up to several parallel rigs at stress rates.
  Metrics: end-to-end latency (median and p95, from source timestamp to browser render), throughput, dropped-sample rate, browser frame rate, and CPU and memory overhead on the source host, compared with a run without the dashboard attached.
  \item[Experiment 2: Filtering fidelity (RQ2).] Spikes, drifts and dropouts of varying amplitude and duration are seeded into traces.
  Moving-average, exponential, median, Hampel and Kalman filters are compared on RMSE against the clean signal, false-alarm rate, missed-event rate and detection delay.
  The result sets the default configuration per channel type.
  \item[Experiment 3: Alert burden (RQ3).] Long replays containing seeded faults are processed with and without the alarm-management rules.
  Metrics: average alerts per 10 minutes, percentage of time in flood, chattering and duplicate alerts, the proportion of seeded critical faults raised at the correct tier, and time-to-alert.
  \item[Experiment 4: Analytical interface evaluation (RQ4).] Representative tasks for each role, such as ``identify the weakest cell'' and ``confirm whether the test passed'', are modelled with KLM in each role view and in a baseline log-and-plot interface, using the system response times measured in Experiment~1.
  Visual clutter is measured from screenshots using feature congestion~\cite{rosenholtz2007clutter}.
  Accessibility is audited with axe-core and a manual WCAG~2.2 AA checklist.
\end{description}

\subsection{Success criteria}\label{subsec:success-criteria}
The targets in Table~\ref{tab:success} are provisional.
They will be confirmed with the STT and the supervisor before the experiments run, and will not be changed after results are seen.

\begin{table}[H]
    \centering\small
    \caption{Provisional success criteria per research question.}
    \label{tab:success}
    \begin{tabularx}{\textwidth}{lY}
        \toprule
        RQ  & Success criterion                                                                                                                                                                                        \\
        \midrule
        RQ1 & p95 end-to-end latency $<$ 500\,ms and zero dropped samples at nominal load; source-host CPU overhead $<$ 5\%; browser $\geq$ 30 frames per second on the engineer view.                                 \\
        RQ2 & The default filter keeps every seeded critical transient visible (missed-event rate 0\%) and reduces RMSE by at least 50\% compared with the raw signal, with detection delay $<$ 1\,s.                  \\
        RQ3 & Average alert rate and time in flood within the ISA-18.2 benchmark levels~\cite{isa2016alarm}; zero missed critical faults; red-tier time-to-alert $<$ 1\,s.                                             \\
        RQ4 & Predicted KLM time for each role's key tasks at least 30\% lower than the baseline; clutter score lower than the baseline; zero critical axe-core violations and a full pass on the manual AA checklist. \\
        \bottomrule
    \end{tabularx}
\end{table}

\todoin{Confirm the ISA-18.2/EEMUA~191 benchmark figures against the standards (via the UL library) and state the exact numbers used for RQ3.}

\subsection{Preliminary results}\label{subsec:preliminary-results}
\todoin{Report any early result honestly, even if small, for example the latency of the two-KPI prototype on one replayed trace, or a first filter comparison on one seeded spike. If no measurement exists yet, say so and state when the first will be taken.}

\subsection{Threats to validity}\label{subsec:threats-to-validity}
\begin{itemize}
    \item \textbf{Construct validity:} KLM predicts expert, error-free performance and cannot measure comprehension or perceived workload.
    A human-subjects study with NASA-TLX~\cite{hart2006nasa} or SAGAT is proposed as future work, noting the critique of NASA-TLX~\cite{babaei2025nasatlx}.
    \item \textbf{Internal validity:} the author builds and evaluates the artefact.
    Mitigated by scripted measurement, pre-registered success criteria and seeded ground truth.
    \item \textbf{External validity:} results on synthetic and public data may not match the STT rig.
    Mitigated by calibrating the simulator's noise and rates to partner traces where permitted.
    \item \textbf{Baseline fairness:} the log-and-plot baseline must represent current practice; it will be built from the STT's existing scripts and checked with the STT.
\end{itemize}
